Esta sección propone una arquitectura para un sistema que, a partir de un repositorio de
\dfn{backend}, produzca un \dfn{frontend} como el de este trabajo. Es una propuesta: \emph{no se
ha implementado ni evaluado} como sistema, y sus componentes se justifican por lo observado
en el proceso de \cref{SEC:GEN-PIPELINE}, no por resultados medidos sobre otros proyectos.

\subsection{Principio de diseño}

La clasificación de \cref{SEC:GEN-CLASSIFICATION} sugiere un principio: \emph{todo lo que
pueda ser determinista debe serlo}, y el razonamiento del modelo se reserva para lo que el
contrato no expresa. Esto no es solo una cuestión de coste. Las etapas deterministas son las
que permiten confiar en el resultado de las demás: el inventario impide que un agente use una
operación que no existe, y la verificación decide sin opinión si un cambio se acepta. En este
trabajo, la regla de no inventar operaciones junto con el inventario generado hizo que el
segundo agente documentara los huecos que encontró en lugar de suplirlos con suposiciones
(\cref{SEC:AN-GAPS}).

\subsection{Componentes}

\paragraph{1. Plantilla de aplicación.} Una base común parametrizable, extraída de la de este
trabajo: cliente de la \acs{api} con renovación de sesión y errores simbólicos
(\texttt{src/lib/api}), enrutador de \dfnpl{mock} (\texttt{src/lib/mock/router.js}),
registro de módulos autodescubiertos (\texttt{src/app/registry.js}), primitivas de interfaz
(\texttt{src/ui}) y auditores (\texttt{scripts/audit}). Los parámetros son pocos y se pueden
leer del contrato y de un número reducido de ficheros del servicio: nombres de los campos de
los \textit{tokens}, \textit{claims} del \acs{jwt}, formato de paginación y formato de errores.

\paragraph{2. Extractor determinista del contrato.} Los \textit{scripts} de
\texttt{scripts/analyze}, generalizados: inventario de operaciones y esquemas, cobertura del
contrato por los módulos y, como extensión natural, generación de las formas de los
\dfnpl{mock} y de tipos a partir de los esquemas.

\paragraph{3. Agente analista del servicio.} Lee el código del \dfn{backend} con una tarea
acotada por el inventario: para cada operación, quién puede invocarla, qué devuelve a cada
actor y qué transiciones provoca. Su salida no es prosa sino datos estructurados (la tabla
actor $\times$ recurso $\times$ acción $\times$ permiso, la máquina de estados y la lista de
huecos), de modo que el resto del sistema la consume sin interpretarla. Cada afirmación
lleva la referencia al fichero del que se dedujo, para que pueda verificarse: en este
trabajo, una supuesta fuga de datos entre organizaciones se descartó precisamente al
releer el código antes de comunicarla (\cref{SEC:AN-BACKEND}).

\paragraph{4. Generador de preguntas.} Convierte los huecos, las reglas ausentes y las
referencias visuales en preguntas cerradas para la persona responsable (\cref{SEC:GEN-CLASSIFICATION}).
Su objetivo es minimizar el número de preguntas, no maximizar la información recogida:
solo pregunta lo que cambia el producto.

\paragraph{5. Planificador.} Construye el grafo entre operaciones y módulos, lo corta en áreas
con pocas dependencias y asigna cada área a un agente, con su lista de rutas propias
(el equivalente a \texttt{OWNERSHIP.json}). La base común se construye antes de repartir.

\paragraph{6. Agentes de implementación.} Uno por área, cada uno en su \dfn{worktree}, con la
guía de estilo como instrucción y el módulo como unidad de trabajo. El repositorio ya contiene
un esbozo de esta idea: las definiciones de subagentes de \texttt{.claude/agents} (analista
del \dfn{backend}, revisor de código, auditor del trabajo paralelo, redactor de
documentación, entre otras) reparten papeles dentro de una misma herramienta, aunque sin el
orquestador que aquí se propone.

\paragraph{7. Oráculo de verificación.} El equivalente a \texttt{npm run check} más la
auditoría de propiedad, como condición necesaria para aceptar cualquier cambio. El registro
de errores muestra que esta condición debe imponerla el entorno y no la disciplina del agente:
en una ocasión se confirmó un cambio con pruebas fallidas porque la orden encadenada no
comprobaba el resultado (entrada 13).

\paragraph{8. Integrador.} Une periódicamente las ramas en una rama temporal y ejecuta el
oráculo sobre la unión, que es donde aparecen los fallos que ninguna rama ve por separado.

\subsection{Límites}

La propuesta hereda los límites del proceso del que se extrae:
\begin{itemize}
  \item \emph{Un solo caso.} Se ha observado un único \dfn{backend} (Django REST con
        drf-spectacular). Servicios sin contrato publicado, con contratos incompletos o con
        la lógica repartida en varios servicios exigirían más razonamiento en las etapas 2 y 4.
  \item \emph{Lo específico no se generaliza.} La anotación sobre el \acs{pdf}, la
        corrección por \dfnpl{rubrica} o la ventana de revisión no salen del contrato: son
        diseño de producto, y aquí ocuparon buena parte del trabajo.
  \item \emph{La estética sigue necesitando a una persona.} El auditor de estilo comprueba
        reglas, no la coherencia visual. Sin capturas de pantalla revisadas, el sistema no
        puede saber si el resultado ``parece'' lo que se pidió.
  \item \emph{Validación real.} Los \dfnpl{mock} se derivan del mismo análisis que el
        código, de modo que un error de interpretación se reproduce en ambos. Sin ejecutar
        contra el servicio real (\cref{SEC:TEST-REAL}), el oráculo solo garantiza coherencia
        interna.
  \item \emph{Orquestación.} El trabajo autónomo de este proyecto se lanzó por iteraciones
        desde un relanzador externo, y en una ocasión llegaron a ejecutarse varias sesiones
        del mismo agente a la vez sobre el mismo \dfn{worktree}, con cambios de una
        arrastrados al commit de otra. Un orquestador real necesita exclusión mutua por
        agente y por directorio de trabajo, algo que el aislamiento por \dfn{worktree} no
        resuelve por sí solo.
\end{itemize}

\pendiente{Valoración del autor: ¿merece la pena prototipar algún componente (por ejemplo, el
analista con salida estructurada) dentro del alcance del TFG o se deja como trabajo futuro?}
